\chapter{Marco Teórico y Metodología}
\label{cap:MarcoTeorico}

% \epigraph{La perfección no se alcanza cuando no hay nada más que añadir, sino cuando no hay nada más que quitar.}{Antoine de Saint-Exupéry}

\lettrine[lraise=-0.1, lines=2, loversize=0.2]{E}{n} este capítulo se establece el marco teórico y metodológico sobre el que se sustenta la arquitectura de detección de esta investigación. Se definen matemáticamente el modelo del sistema y la formulación del canal de radiofrecuencia, estableciendo los retos inherentes al régimen de baja relación señal a ruido (\textit{\gls{snr}}). A continuación se detalla el conjunto de datos \textit{NoisyUAV v2}, analizando su estructura y sus características. Sobre esta base, el capítulo desarrolla tres pilares metodológicos: la justificación del procesamiento directo de secuencias complejas I/Q frente a espectrogramas, el modelado del ruido de etiquetas dependiente de la dispersión temporal, y la formulación teórica del preprocesamiento basado en entropía de Shannon y de las Redes Neuronales Convolucionales de Valores Complejos (\textit{\gls{cv-cnn}}).

\section{Modelo del Sistema y Formulación del Canal RF}
\label{sec:modelo_sistema}
El diseño de un sistema robusto de detección pasiva requiere una caracterización rigurosa del entorno electromagnético y de las señales involucradas. En este proyecto, se asume un escenario de monitorización donde un receptor de radio definida por \textit{software} \gls{sdr} escanea de forma continua la banda {\gls{ism}} de 2.4\,GHz.

En un escenario de monitorización no colaborativo, la señal captada por la antena receptora en el instante $t$, denotada como $r(t)$, se modela analíticamente como la superposición de la señal de interés (sometida a la distorsión del canal de propagación), las interferencias cocanal coexistentes y el ruido térmico del receptor \cite{rout_novel_2025, ezuma_detection_2020}. En su representación de banda base equivalente, el sistema se define mediante la ecuación:

\begin{equation}
 r(t) = \big(h * x\big)(t) + i(t) + n(t)    
 \label{eq:modelo_sistema}
\end{equation}

Donde cada término representa una magnitud física discreta del canal inalámbrico:
\begin{itemize}
    \item $x(t)$ es la señal compleja en banda base transmitida por el dron o su radiocontrol. Típicamente, estas señales emplean esquemas \gls{fhss} acoplados sobre modulaciones digitales de capa física de envolvente constante, sobre todo se usan modulación por desplazamiento de frecuencia gaussiana (gls{gfsk}) y modulación por desplazamiento de frecuencia ({fsk}), o derivados como LoRa; con el objetivo de mitigar interferencias y evitar intercepciones, tal y como constatan los análisis de \gls{rf} en drones comerciales realizados en \cite{ezuma_detection_2020, gluge_robust_2024}.
    \item $h(t)$ representa la respuesta al impulso compleja del canal físico en banda base equivalente. Este término compacta los fenómenos de atenuación por propagación (\textit{path loss}), rotación fasorial de fase y desvanecimiento destructivo/constructivo intersímbolo, modelando las transformaciones que sufre la envolvente compleja de la señal antes de la etapa de digitalización en el receptor.
    \item $i(t)$ representa la interferencia co-canal. Este término agrupa las emisiones de otros dispositivos coexistentes en la banda de interés, principalmente ráfagas de redes IEEE 802.11 (Wi-Fi) y Bluetooth \cite{ezuma_detection_2020}.
    \item $n(t)$ modela el ruido térmico inherente a la electrónica del receptor \cite{lacy_machine_2024}.
\end{itemize}

Dado que la Ecuación~\ref{eq:modelo_sistema} se formula en el dominio de banda base equivalente, es preciso realizar una caracterización estadística rigurosa del término $n(t)$. Como se detalla en la literatura de referencia sobre teoría de detección \cite{kay_fundamentals_1998}, un proceso aleatorio gaussiano complejo requiere para su descripción completa no solo de su media y su covarianza, sino también de su pseudocovarianza. En las arquitecturas de recepción de radiofrecuencia, el ruido térmico captado por la antena se traslada a banda base mediante un proceso de demodulación en cuadratura. Esta conversión divide la señal en dos ramas ortogonales (generalmente mediante la multiplicación por portadoras desfasadas noventa grados), lo que garantiza que las componentes resultantes en fase ($I$) y en cuadratura ($Q$) sean estadísticamente independientes y posean la misma varianza. Como resultado, la pseudocovarianza del proceso se anula, definiendo a $n(t)$ como un ruido gaussiano complejo propio o circularmente simétrico (\gls{cscg}). De manera formal, este proceso se denota como $n(t) \sim \mathcal{CN}(0, \sigma_n^2)$, indicando una media nula y una potencia total $\sigma_n^2$, donde la potencia de las componentes real e imaginaria corresponde a $\sigma_n^2/2$ en cada una de ellas.

Para su procesamiento digital en las arquitecturas de aprendizaje profundo, la señal analógica $r(t)$ es muestreada a una frecuencia $F_s$, obteniendo la secuencia discreta $r[n]$. Para preservar la coherencia de fase relativa y aislar las firmas electromagnéticas del transmisor (\textit{hardware impairments}), la señal se procesa directamente como una secuencia de valores complejos compuesta por sus componentes en fase ($r_I$) y en cuadratura ($r_Q$) \cite{valkama_blind_2005, zheng_deep_2025}:
\begin{equation}
    r[n] = r_I[n] + j r_Q[n]
    \label{eq:iq_discreta}
\end{equation}

\subsection{Definición del régimen de baja Relación Señal a Ruido (\gls{snr})}
El desafío fundamental de este proyecto radica en la detección de la ráfaga $x(t)$ en escenarios donde su potencia puede ser inferior a la de las transmisiones interferentes. La Relación Señal a Ruido e Interferencia (\gls{sinr}), que en el contexto de esta investigación y en la bibliografía sobre detección de drones se denota de forma generalizada simplemente como \gls{snr}, se formula como:
\begin{equation}
    \text{SNR (dB)} = 10 \log_{10} \left( \frac{P_x}{P_i + P_n} \right)
    \label{eq:snr_calculo}
\end{equation}

Donde los escalares $P_x$, $P_i$ y $P_n$ representan la potencia media total (obtenida como la integral de la \gls{psd} sobre el ancho de banda del receptor) de la señal de interés, la interferencia cocanal y el ruido térmico, respectivamente. 

En un escenario de recepción ideal, la presencia de la señal incrementa la potencia total medida, permitiendo el uso de detectores de energía clásicos ya mencionados en este trabajo. Sin embargo, cuando la \gls{snr} desciende por debajo de los 0\,dB, la variable de decisión del receptor queda matemáticamente dominada por la varianza del ruido $P_n$ y de la interferencia $P_i$. En este régimen de \gls{snr} negativa, la envolvente de la señal discreta $r[n]$ resulta indistinguible de una distribución gaussiana, lo que provoca que los detectores clásicos colapsen y generen tasas inasumibles de falsas alarmas \cite{lacy_machine_2024}.

\section{El conjunto de datos \textit{NoisyUAV v2}}
\label{sec:dataset_noisyuav}

Tras el análisis crítico de las alternativas, se ha seleccionado el conjunto de datos \textit{NoisyUAV v2} como el eje experimental de este proyecto. Esta base de datos se diseñó con el objetivo explícito de evaluar la robustez algorítmica en entornos de muy baja \gls{snr} y con interferencias reales. Se resumen las principales características:

Autores y artículo: Glüge et al. (2024) \cite{gluge_robust_2024}.

Arquitectura de adquisición de RF: captación realizada empleando una plataforma \gls{sdr} basada en componentes comerciales, concretamente el receptor Ettus Research USRP B210. Las señales se capturaron con una frecuencia de muestreo de 56\,MHz, correspondiente al máximo ancho de banda para capturas continuas en tiempo real con el que cuenta este dispositivo. Posteriormente, con el objetivo explícito de reducir el consumo de memoria y la carga computacional en las fases posteriores de entrenamiento, los autores redujeron este valor mediante submuestreo desde 56\,MHz hasta 14\,MHz, aplicando previamente un filtro anti-aliasing Chebyshev de tipo\,I y orden\,8. La frecuencia de muestreo efectiva tras este procesamiento es, por tanto, $F_s = 14$\,MHz \cite{gluge_robust_2024}.

Volumen, topología y formato de los datos: el conjunto comprende 17.744 grabaciones almacenadas como vectores IQ. Cada archivo representa una ventana de observación continua de $2^{20} = 1.048.576$ muestras obtenidas a una frecuencia de 14\,MHz, lo que equivale a una ventana de observación de aproximadamente 74,9\,ms. El conjunto de datos contiene transmisiones de seis modelos de drones comerciales de distinta gama y marca como DJI, Futaba, Graupner, Taranis o Turnigy; todos ellos empleando distintas variaciones protocolos propietarias basados en \gls{fhss}. A lo largo del documento, se referenciará a estos seis tipos de transmisores como Dron o \textit{Target} 0, 1, 2, 3, 5 y 6 indistintamente. También cuenta con una clase que únicamente incluye grabaciones en las que no se producen transmisiones de drones, es decir, sólo se registraron interferencias y ruido del canal (Wi-Fi, Bluetooth y \gls{awgn}), se trata del \textit{Target} 4. No obstante, para los fines de esta investigación, el problema se abordará estrictamente desde una perspectiva de detección binaria donde los seis transmisores se agruparán en un único grupo, la Clase 1: presencia de dron; mientras que todas las grabaciones del \textit{Target} 4 y la no detección de dron en las grabaciones de cualquier otro \textit{Target} constituirán la Clase 0:ausencia de dron.

Análisis de la hostilidad electromagnética: tal y como se ha detallado en el Capítulo~\ref{cap:EstadoDelArte}, el rendimiento de los sistemas de detección RF presenta una degradación característica en función del nivel de \gls{snr}, con umbrales de colapso que la literatura sitúa de forma consistente en torno a los $0$\,dB para detectores clásicos y entre $-10$ y $-15$\,dB para aproximaciones basadas en aprendizaje profundo \cite{zheng_uavsignal_2025, gluge_robust_2024, vasant_ahirrao_rf-based_2024}. \textit{NoisyUAV v2} cubre íntegramente este espectro, desde condiciones favorables hasta regímenes extremos, permitiendo una validación exhaustiva organizada en cuatro grupos:

\begin{itemize}
    \item Grupo A (\gls{snr} $\ge$ 6\,dB): el caso óptimo, en tanto que la potencia es significativamente superior al suelo de ruido. Sin embargo, las interferencias presentan perfiles energéticos equiparables a las transmisiones de los drones, con características muy similares en el caso de los \textit{bursts} de Bluetooth.
    \item Grupo B ($0$\,dB $\le$ \gls{snr} $<$ 6\,dB): complejidad moderada. En este rango de \gls{snr}, la coexistencia co-canal degrada la integridad fasorial de forma mucho más agresiva debido a que la señal de interés se encuentra muy enmascarada por el ruido de fondo y las interferencias existentes. Este rango de valores se considera la línea base para sistemas comerciales que realizan detecciones coherentes\footnote{En telecomunicaciones, un receptor coherente es aquel que realiza un seguimiento activo y recupera la fase de la portadora de la señal transmitida, utilizando esta referencia de fase exacta para el proceso de demodulación.} \cite{vasant_ahirrao_rf-based_2024, martelli_exploitation_2020}.
    \item Grupo C (-12\,dB $\le$ \gls{snr} $<$ 0\,dB): en este entorno, los detectores clásicos basados en energía son incapaces de operar correctamente. Esta es la razón por la que el diseño de arquitecturas de \textit{Deep Learning} se centra en este segmento de \gls{snr}. Como se vio en el Capítulo \ref{cap:EstadoDelArte}, la literatura predominante emplea \glspl{cnn}, diseñadas para visión artificial, sobre espectrogramas. En este trabajo se emplearán arquitecturas de \textit{Deep Learning} orientadas a extraer características directamente de los tensores IQ, evitando la dependencia de estos modelos de imagen, constituyendo este grupo el eje de evaluación central \cite{gluge_robust_2024}.
    \item Grupo D (\gls{snr} $<$ -12\,dB): las detecciones comienzan a ser marginales, ocurriendo con muy baja reproducibilidad \cite{gluge_robust_2024}. La potencia de la señal es estadísticamente inferior a la varianza del ruido del canal, dificultando la extracción de patrones  y características que contengan suficiente información. Este grupo se utilizará para evaluar el umbral límite de robustez del sistema diseñado.
\end{itemize}

Justificación de la selección: constituye el único conjunto de datos que satisface íntegramente los requisitos definidos. Los 14\,MHz de $F_s$ evitan la redundancia computacional, los 75\,ms continuos permiten el escaneo activo, y la presencia inherente de regímenes de \gls{snr} negativa garantiza una validación robusta de las arquitecturas propuestas.

\subsection{Inspección del conjunto de datos seleccionado}

Para comprender el escenario de detección, resulta fundamental analizar la naturaleza fasorial y espectral de las señales. La Tabla \ref{tab:distribucion_muestras} desglosa el número de muestras, evidenciando un conjunto de datos perfectamente balanceado para el problema de clasificación binaria propuesto. Adicionalmente, la tabla especifica la frecuencia central ($f_c$), el ancho de banda ($\Delta f$) y la duración media temporal de las ráfagas (\textit{bursts}) emitidas por cada modelo de radiocontrol. Cabe destacar que la clase correspondiente a ruido e interferencias (correspondiente al Identificador 4) omite estos parámetros.

\begin{table}[htbp]
\centering
\small
\caption{Distribución de muestras en el conjunto de datos \textit{NoisyUAV v2}.}
\label{tab:distribucion_muestras}
\resizebox{\linewidth}{!}{%
\begin{tabular}{lccccc}
\toprule
\textbf{Origen} & \textbf{Identificador clase} & \textbf{Muestras} & \textbf{$f_c$ (GHz)} & \textbf{$\Delta f$ (MHz)} & \textbf{Duración (ms)} \\ \midrule
Dron DJI Phantom GL300F Pro & 0 & 1280 & 2.44175 & 1.2 & 2.18\\
Dron Futaba T7C & 1 & 3472 & 2.44175 & 2.0 & 1.70\\
Dron Futaba T14SG & 2 & 801 & 2.44175 & 3.1 & 1.40\\
Graupner mx-16 & 3 & 801 & 2.44175 & 1.0 & 1.9/3.7\\
Ruido & 4 & 8872 & - & - & - \\
Dron Taranis ACCST & 5 & 1663 & 2.44000 & 1.5 & 3.1/4.4\\
Dron Turnigy 9X & 6 & 855 & 2.44500 & 2.0 & 1.30\\ \bottomrule
\textbf{Total} & & \textbf{17.744} & & & \\ \bottomrule
\end{tabular}
}
\end{table}

\textbf{1. Variabilidad interclase y regímenes de operación:} un sistema de detección debe ser capaz de identificar las transmisiones sin una dependencia estricta de la potencia de recepción. En condiciones de alta relación señal a ruido, las propiedades de la secuencia temporal permiten discernir el modelo específico de radiocontrol al que pertenece la ráfaga, incluso analizando directamente su amplitud. Este escenario óptimo permite que el sistema basado en \gls{dl} extraiga las características comunes subyacentes a los protocolos propietarios de los seis transmisores, aprendiendo simultáneamente a discriminar dichas transmisiones frente al ruido y las interferencias. Sin embargo, como muestra la Figura \ref{fig:comparativa_clases_completa}, la reducción de la \gls{snr} distorsiona tanto la señal temporal como sus propiedades frecuenciales (visibles en los periodogramas y espectrogramas). En consecuencia, a bajos niveles de \gls{snr}, las diferencias entre los distintos transmisores, así como su separabilidad frente a las interferencias o al ruido del canal, se reducen drásticamente. 

\enlargethispage{1.5cm} % TRUCO 1: Alargamos el folio virtualmente 2 cm para que quepa la figura ampliada

\begin{figure}[!htbp]
    \centering
    \begin{subfigure}[b]{0.49\linewidth} % TRUCO 2: Aumentamos el ancho
        \centering
        \includegraphics[width=\linewidth]{figuras/seccion3/comparacion_clase_28dB.png}
        \caption{Régimen favorable (28\,dB).}
    \end{subfigure}
    \hfill
    \begin{subfigure}[b]{0.49\linewidth}
        \centering
        \includegraphics[width=\linewidth]{figuras/seccion3/comparacion_clase_10dB.png}
        \caption{Régimen intermedio (10\,dB).}
    \end{subfigure}

    \vspace{0.3cm} 

    \begin{subfigure}[b]{0.49\linewidth}
        \centering
        \includegraphics[width=\linewidth]{figuras/seccion3/comparacion_clase_0dB.png}
        \caption{Régimen moderado (0\,dB).}
    \end{subfigure}
    \hfill
    \begin{subfigure}[b]{0.49\linewidth}
        \centering
        \includegraphics[width=\linewidth]{figuras/seccion3/comparacion_clase_-12dB.png}
        \caption{Régimen crítico ($-12$\,dB).}
    \end{subfigure}
    
    \caption{Comparación por clase en función de la degradación progresiva de la relación señal a ruido.}
    \label{fig:comparativa_clases_completa}
\end{figure}

\textbf{2. Análisis multidominio de la degradación:} para ilustrar el impacto del ruido térmico y las interferencias sobre las transmisiones, la Figura~\ref{fig:degradacion_target0_total} expone la evolución de las transmisiones del \textit{Target~0} (correspondiente al dron DJI Phantom GL300F Pro, véase la Tabla~\ref{tab:distribucion_muestras}) en los dominios complejo I/Q, temporal (envolvente) y tiempo-frecuencia (espectrograma y \gls{psd}). Esta transición muestra la degradación progresiva de la señal a medida que la \gls{snr} disminuye. En el dominio temporal, a valores altos de SNR, la ubicación de las ráfagas es directa; sin embargo, en regímenes críticos, la envolvente pasa a estar dominada por la varianza del ruido de fondo. Por su parte, la representación tiempo-frecuencia (espectrograma) exhibe una mayor tolerancia matemática a la disminución de la \gls{snr}, preservando la estructura de las ráfagas. Esta característica explica que la mayor parte de las metodologías del estado del arte transformen la señal de \gls{rf} en imágenes espectrales para procesarlas con redes orientadas a visión artificial.

\begin{figure}[!htbp]
    \centering
    \begin{subfigure}[b]{0.495\linewidth} % Apuramos al máximo el ancho de la columna
        \centering
        \includegraphics[width=\linewidth]{figuras/seccion3/ej_dataset_target0_30db.png}
        \caption{\textit{Target 0} a 30\,dB.}
    \end{subfigure}
    \hfill
    \begin{subfigure}[b]{0.495\linewidth}
        \centering
        \includegraphics[width=\linewidth]{figuras/seccion3/ej_dataset_target0_0db.png}
        \caption{\textit{Target 0} a 0\,dB.}
    \end{subfigure}

    \vspace{0.3cm} % Reducimos un poco el espacio central

    \begin{subfigure}[b]{0.495\linewidth}
        \centering
        \includegraphics[width=\linewidth]{figuras/seccion3/ej_dataset_target0_-12db.png}
        \caption{\textit{Target 0} a $-12$\,dB.} % He corregido el signo menos matemático aquí
    \end{subfigure}
    \hfill
    \begin{subfigure}[b]{0.495\linewidth}
        \centering
        \includegraphics[width=\linewidth]{figuras/seccion3/ej_dataset_target0_-20db.png}
        \caption{\textit{Target 0} a $-20$\,dB.} % He ajustado el texto al nombre del archivo (-20dB)
    \end{subfigure}
    \caption{Evolución multidominio de la clase objetivo 0 en función del nivel de ruido del canal.}
    \label{fig:degradacion_target0_total}
\end{figure}

% \clearpage % Este comando es un muro absoluto: empuja todo lo que venga debajo a la página siguiente

\textbf{3. Caracterización de las interferencias y falsos positivos:} la Figura \ref{fig:degradacion_ruido_total} detalla el comportamiento multidominio de la clase correspondiente a ruido e interferencias. Como se ha detallado anteriormente, a altos niveles de \gls{snr} (Figura \ref{fig:degradacion_ruido_total}a y \ref{fig:degradacion_ruido_total}b), las transmisiones Bluetooth tienen transmisiones análogas a las ráfagas \gls{fhss} de los drones, mientras que las transmisiones Wi-Fi inyectan pulsos capaces de saturar la recepción y provocar falsos positivos. En contrapartida, en regímenes de baja \gls{snr} (Figura \ref{fig:degradacion_ruido_total}d), la potencia del ruido \gls{awgn} supera al resto de transmisiones, dificultando la distinción entre la señal de interés y las interferencias cocanal en cualquier dominio.
\begin{figure}[!htbp]
    \centering
    \begin{subfigure}[b]{0.495\linewidth}
        \centering
        \includegraphics[width=\linewidth]{figuras/seccion3/ej_dataset_target4_30db.png}
        \caption{Clase ruido a 30\,dB.}
    \end{subfigure}
    \hfill
    \begin{subfigure}[b]{0.495\linewidth}
        \centering
        \includegraphics[width=\linewidth]{figuras/seccion3/ej_dataset_target4_12db.png}
        \caption{Clase ruido a 12\,dB.}
    \end{subfigure}

    \vspace{0.3cm} % Espacio reducido para compensar el aumento de altura

    \begin{subfigure}[b]{0.495\linewidth}
        \centering
        \includegraphics[width=\linewidth]{figuras/seccion3/ej_dataset_target4_0db.png}
        \caption{Clase ruido a 0\,dB.}
    \end{subfigure}
    \hfill
    \begin{subfigure}[b]{0.495\linewidth}
        \centering
        \includegraphics[width=\linewidth]{figuras/seccion3/ej_dataset_target4_-12db.png}
        \caption{Clase ruido a $-12$\,dB.}
    \end{subfigure}
    \caption{Análisis de las emisiones interferentes y ruido de fondo a distintos regímenes de potencia.}
    \label{fig:degradacion_ruido_total}
\end{figure}

\section{Justificación del procesamiento directo de la señal IQ  frente a representaciones de tiempo-frecuencia}
\label{sec:justificacion_procesamiento_IQ}

El paradigma predominante en la literatura reciente orientada a la clasificación de señales electromagnéticas mediante aprendizaje profundo, tal y como se detalla en el Capítulo \ref{cap:EstadoDelArte}, se fundamenta en la proyección de la serie temporal unidimensional hacia una representación bidimensional en el dominio tiempo-frecuencia. En este contexto, el estimador de densidad espectral más empleado es el espectrograma \cite{gluge_robust_2024, zhao_drone_2024}. Matemáticamente, es una matriz que se calcula evaluando la magnitud al cuadrado de la \gls{stft}. Esta constituye una extensión de la \gls{dtft} que permite analizar señales no estacionarias cuyas propiedades espectrales varían en el tiempo. La idea fundamental consiste en segmentar la señal en fragmentos temporales de una longitud determinada mediante una función de ventana $w[n]$ (Rectangular, Hanning, Hamming y Blackman son las más usadas \cite{zhang_window, harris_window}), aplicar la transformada a cada segmento de forma independiente y organizar los espectros resultantes en una matriz bidimensional. Formalmente, para una señal discreta $r[n]$, su espectrograma $S[m,k]$ se define como:
\begin{equation}
    S[m,k] = \left| \sum_{n=0}^{N_w-1} r[n+mH] w[n] e^{-j\frac{2\pi}{N_w}kn} \right|^2
    \label{eq:espectrograma}
\end{equation}
En esta ecuación, $m$ representa el índice de la trama temporal y $k$ es el \textit{bin}\footnote{Un \textit{bin} de frecuencia representa cada uno de los intervalos discretos en los que se divide el espectro continuo de la señal al aplicar la \gls{dtft}. Su anchura ($\Delta f$) determina la resolución frecuencial, es decir, el tamaño mínimo en frecuencia que tiene el sistema.} de frecuencia. El término $w[n]$ denota la función de la ventana aplicada, $N_w$ determina la longitud de dicha ventana, y el parámetro $H$ rige el desplazamiento temporal discreto entre segmentos consecutivos. Este último factor controla la tasa de solapamiento de las ventanas, dictando la resolución temporal del estimador frente a la evolución transitoria de la señal. La resolución combinada de esta matriz queda estrictamente sujeta al principio de incertidumbre de Heisenberg-Gabor: ventanas largas ofrecen alta resolución frecuencial pero baja resolución temporal, y viceversa \cite{oppenheim_discrete_2014}.

La motivación teórica que subyace a esta transformación espacial radica en la posibilidad de tratar la matriz espectral resultante como una imagen digital fotométrica. Este isomorfismo permite importar e implementar arquitecturas de redes neuronales convolucionales que fueron concebidas para tareas de reconocimiento de imágenes. En dicho marco teórico, el propósito de las capas de convolución es recorrer la matriz espacial bidimensional para identificar jerarquías de características topológicas, asimilando las variaciones de densidad espectral a gradientes espaciales, bordes o patrones de textura visual. 

No obstante, la asimilación conceptual de una señal de radiofrecuencia a un mapa de píxeles carece de rigor físico, induciendo deficiencias estructurales en el modelo. Concretamente, la adopción del espectrograma como bloque de entrada impone tres limitaciones fundamentales que comprometen severamente la detección en regímenes donde la relación señal a ruido desciende por debajo de los 0\,dB:

\begin{itemize}
    \item Destrucción de la coherencia de fase: La operación del valor absoluto en la Ecuación \ref{eq:espectrograma} descarta la fase instantánea de la señal. En regímenes críticos ($\text{SNR} \le -6$\,dB), las transiciones de fase inherentes a la modulación y las firmas de \textit{hardware} de los osciladores locales (\textit{RF fingerprinting}) constituyen discriminantes estadísticamente vitales frente al ruido gaussiano que resultan irrecuperables tras la transformación espectral \cite{zheng_deep_2025}.

    \item Límite de Heisenberg-Gabor: La \gls{stft} está sujeta al principio de incertidumbre tiempo-frecuencia, que establece un compromiso ineludible: una alta resolución frecuencial (ventanas largas) diluye los bordes transitorios de las ráfagas \gls{fhss}, mientras que una alta resolución temporal (ventanas cortas) ensancha artificialmente el espectro, incrementando el solapamiento con interferencias co-canal \cite{oppenheim_discrete_2014}.

    \item Complejidad espacial y carga computacional: el procesamiento basado en espectrogramas incrementa el coste computacional en dos etapas. Por un lado, la obtención de la \gls{stft} requiere una complejidad de $\mathcal{O}(L \log L)$ para una secuencia de longitud $L$. Por otro, el posterior procesado mediante \gls{cnn} 2D presenta una carga computacional que escala de forma cuadrática, $\mathcal{O}(K^2)$, respecto al tamaño del \textit{kernel} espacial \cite{alam_rf-enabled_2023}. En contraste, operar directamente sobre las secuencias IQ en banda base permite el uso de redes convolucionales unidimensionales, eliminando la necesidad de calcular transformadas previas y reduciendo la complejidad de las convoluciones a un orden lineal $\mathcal{O}(K)$ \cite{zheng_deep_2025}. Estas características producen que este tipo de modelos puedan desplegarse en nodos \textit{Edge AI}, donde los recursos de hardware y potencia son estrictamente limitados \cite{cai_toward_2024, zhang_edge_2023}.
\end{itemize}

Pese a las ventajas de mantener la información compleja de la señal y reducir la carga de trabajo en el entrenamiento y la inferencia, operar únicamente en el dominio del tiempo puede dificultar la detección cuando existen muchas interferencias o la \gls{snr} es muy baja. Por ello, la literatura reciente propone el uso de arquitecturas duales o multidominio \cite{liu_method_2026}. Estas redes combinan una rama principal para procesar la señal temporal con una rama secundaria que extrae características a partir de otras representaciones de la señal, como la \gls{psd} (enfoque que se detallará en el Capítulo \ref{cap:arquitectura}).

\section{Ruido en el etiquetado de las transmisiones de dron}
\label{sec:label_noise}

Tal y como se ha descrito en la Sección~\ref{sec:dataset_noisyuav}, cada grabación de \textit{NoisyUAV v2} tiene una duración de $T_{\text{obs}} = 75$\,ms ($2^{20} = 1.048.576$ muestras a $F_s = 14$\,MHz) e incluye, además de las transmisiones del dron, ruido térmico e interferencias cocanal de redes Wi-Fi y Bluetooth. El problema que se analiza en esta sección surge porque las transmisiones de dron ocupan únicamente una fracción reducida de esa grabación. Como se cuantifica a continuación, incluso en el caso más favorable, el enlace de control del dron está activo durante menos del 12\,\% del total de muestras. El 88\,\% restante contiene exclusivamente ruido térmico e interferencias cocanal. Si el modelo recibe como entrada la secuencia completa de 75\,ms, la función de pérdida durante el entrenamiento queda dominada por los tramos sin transmisión de dron, que constituyen la mayor parte del tensor. Al asignar la etiqueta de clase Dron (clase 1) a una grabación cuya mayor parte no contiene señal de dron, el optimizador asocia patrones de ruido e interferencia con la clase objetivo. Este fenómeno se denomina en la literatura del aprendizaje máquina como ruido de etiquetas dependiente de la instancia (\textit{instance-dependent label noise}) \cite{frenay_classification_2014, song_learning_2022}.

Para cuantificar la ocupación temporal de las transmisiones de dron dentro de cada grabación del conjunto de datos seleccionado, se define la ocupación máxima como el cociente entre el tiempo total de transmisión activa y la duración de la grabación:

\begin{equation}
    \text{Ocupación máxima (\%)} = \frac{N_{\text{burst}} \cdot \max(T_{\text{burst}})}{T_{\text{obs}}} \times 100
    \label{eq:ocupacion_maxima}
\end{equation}

Donde $N_{\text{burst}}$ es el número estimado de transmisiones de dron presentes en una grabación de $T_{\text{obs}} = 75$\,ms, y $\max(T_{\text{burst}})$ es la duración máxima de cada transmisión de dron. Aplicando la Ecuación~\ref{eq:ocupacion_maxima} a los parámetros de la Tabla~\ref{tab:distribucion_muestras}, se obtiene la ocupación temporal para cada modelo del conjunto de datos, tal y como se detalla en la Tabla~\ref{tab:ocupacion_temporal}.

\begin{table}[htbp]
\centering
\small
\caption{Ciclo de trabajo empírico y ocupación temporal máxima por modelo en una ventana de observación de 75\,ms.}
\label{tab:ocupacion_temporal}
\resizebox{0.9\linewidth}{!}{%
\begin{tabular}{lccc}
\toprule
\textbf{Dron / Controlador} & \textbf{Ráfagas estimadas ($N_{\text{burst}}$)} & \textbf{Duración máxima $\max(T_{\text{burst}})$ (ms)} & \textbf{Ocupación máxima (\%)} \\ \midrule
Turnigy 9X & 6 & 1.30 & 10.40\% \\
Taranis ACCST & 2 & 4.40 & 11.73\% \\
Graupner mx-16 & 2 & 3.70 & 9.86\% \\
Futaba T14SG & 2 & 3.10 & 8.26\% \\
DJI Phantom GL300F & 2 & 2.18 & 5.81\% \\
Futaba T7C & 2 & 2.00 & 5.33\% \\ \bottomrule
\end{tabular}
}
\end{table}

La Tabla~\ref{tab:ocupacion_temporal} muestra que la ocupación máxima corresponde al modelo Taranis ACCST, con un 11,73\,\%, debido a que sus transmisiones individuales tienen una duración mayor que el resto de modelos del conjunto de datos. El modelo Turnigy 9X alcanza una ocupación del 10,40\,\% con 6 transmisiones por ventana, a pesar de que cada una dura solo 1,30\,ms. Por tanto, en el caso más favorable, el dron ocupa como máximo el 11,73\,\% del tiempo analizado. El 88,27\,\% restante, aproximadamente $925.574$ muestras, contiene únicamente ruido térmico e interferencias de Wi-Fi y Bluetooth, tal y como se caracterizó en la Sección~\ref{sec:dataset_noisyuav}.

En \textit{NoisyUAV v2}, a cada grabación de 75\,ms le corresponde una única etiqueta binaria $\tilde{y} \in \{0,1\}$, donde $\tilde{y} = 1$ indica presencia de dron (clase 1) y $\tilde{y} = 0$ indica ausencia de dron (clase 0) \cite{gluge_robust_2024}. Esa etiqueta se aplica a la totalidad de las muestras de la ventana, sin distinción entre los tramos con transmisión activa del dron y los tramos de ruido e interferencia. El efecto de este esquema de etiquetado depende directamente del nivel de SNR. Cuando la SNR es alta, la potencia de la señal del dron $P_x$ supera a la de las interferencias $P_i$ y el ruido térmico $P_n$ (definidos en la Ecuación~\ref{eq:snr_calculo}), por lo que las características de la señal del dron son distinguibles en el tensor de entrada. Sin embargo, cuando $\text{SNR} < 0$\,dB, se cumple que $P_x < P_i + P_n$: la potencia de la señal del dron es inferior al nivel conjunto de ruido e interferencias, y la etiqueta $\tilde{y} = 1$ se asigna a grabaciones cuya mayor parte no contiene señal de dron detectable.

Formalmente, esto implica que la distribución de probabilidad de la etiqueta $\tilde{y}$ diverge de la clase electromagnética realmente dominante en cada tramo de la grabación:

\begin{equation}
    P(\tilde{y} = k \mid \mathbf{x}, y = c) \neq \delta_{kc}
    \label{eq:label_noise}
\end{equation}

Donde $k$ es la etiqueta asignada a la grabación, $c$ es la clase electromagnética realmente dominante en el tramo evaluado, $\mathbf{x} \in \mathbb{C}^N$ es el vector de muestras de entrada, y $\delta_{kc}$ es la delta de Kronecker.

El impacto del ruido de etiquetas en el entrenamiento se formaliza mediante el principio de Minimización del Riesgo Empírico (\textit{Empirical Risk Minimization}, \gls{erm}) \cite{vapnik_nature_1998}. Este término fundamental del aprendizaje máquina establece que, dado que un algoritmo no puede conocer la distribución estadística completa de los datos, debe aprender minimizando el error sobre el conjunto finito de datos de entrenamiento disponible. Bajo este paradigma, el objetivo teórico del optimizador es minimizar el riesgo esperado $R(\theta)$ sobre la distribución conjunta real $P(\mathbf{x}, y)$:

\begin{equation}
    R(\theta) = \mathbb{E}_{P(\mathbf{x}, y)} \left[ \mathcal{L}(f_\theta(\mathbf{x}), y) \right]
    \label{eq:riesgo_teorico}
\end{equation}

Donde $\mathcal{L}$ es la función de pérdida (por ejemplo, entropía cruzada) y $f_\theta(\mathbf{x})$ es la salida del modelo. Sin embargo, el optimizador no tiene acceso a $P(\mathbf{x}, y)$, sino a la distribución empírica $\tilde{P}(\mathbf{x}, \tilde{y})$ afectada por el ruido de etiquetas. Por tanto, minimiza el riesgo empírico evaluado sobre estas etiquetas, el cual se denota como $\tilde{R}(\theta)$ para distinguirlo del riesgo teórico:

\begin{equation}
    \tilde{R}(\theta) = \frac{1}{M} \sum_{i=1}^{M} \mathcal{L}(f_\theta(\mathbf{x}_i), \tilde{y}_i)
    \label{eq:riesgo_empirico}
\end{equation}

Donde $M$ es el número de grabaciones de entrenamiento y $\tilde{y}_i$ es la etiqueta asignada a la grabación $\mathbf{x}_i$. Dado el ruido de etiquetas formulado en la Ecuación~\ref{eq:label_noise}, se demuestra que el riesgo empírico es un estimador sesgado del riesgo verdadero: $\mathbb{E}[\tilde{R}(\theta)] \neq R(\theta)$ \cite{natarajan_learning_2013, song_learning_2022}.

La consecuencia directa de este sesgo es que, cuando el modelo recibe repetidamente grabaciones de 75\,ms en las que el ruido del canal o las transmisiones de Wi-Fi, Bluetooth son las señales dominantes pero la etiqueta asignada es $\tilde{y} = 1$ (clase Dron), la minimización de $\tilde{R}(\theta)$ lleva al clasificador a aprender los patrones del \gls{awgn} o de las interferencias como características de la clase Dron. Este fenómeno se denomina aprendizaje basado en atajos \textit{shortcut learning}: el modelo converge a un mínimo de la función de pérdida que no corresponde a las características propias de las transmisiones del dron, sino a las de las señales que dominan la mayor parte de las muestras de entrenamiento \cite{geirhos_shortcut_2020}.

La eliminación del ruido de etiquetas requeriría recortar manualmente las transmisiones de dron en cada grabación del conjunto de datos, tarea que para el conjunto de datos seleccionado tiene gran dificultad debido al gran número de grabaciones que contiene. Como alternativa automatizada, el sistema propuesto incorpora un detector analítico de ráfagas como primera etapa. Su función es aislar las transmisiones activas y descartar los tramos que contienen únicamente ruido. Este detector no clasifica la señal (no distingue entre dron e interferencia cocanal), sino que delega esta tarea al modelo que realiza la detección. Su objetivo es garantizar que la red neuronal procese únicamente ráfagas electromagnéticas y no ventanas constituidas en su totalidad por \gls{awgn}. Este mecanismo se detalla en la Sección~\ref{sec:detector_entropia}.

\section{Detección transmisiones mediante Entropía Espectral de Shannon}
\label{sec:detector_entropia}

Tal y como se ha formalizado en la Sección~\ref{sec:label_noise}, la detección en el límite del ruido térmico exige un mecanismo de segmentación temporal previo al clasificador. Este preprocesador no evalúa la naturaleza del emisor, sino que rastrea la serie temporal para aislar las regiones de interés con concentración espectral anómala. En el contexto de los protocolos \gls{fhss}, este requisito puede satisfacerse de forma analítica explotando una propiedad fundamental de la información de Shannon \cite{zhang_entropy-based_1989}: la entropía espectral de la señal observada cae bruscamente en el instante exacto de una transmisión activa.

A diferencia de los detectores de energía convencionales, basados en la comparación de la potencia total integrada frente a un umbral fijo, el enfoque propuesto opera sobre la distribución de la energía en el dominio de la frecuencia. Este principio ha sido aplicado en trabajos previos sobre análisis espectral de señales de drones y \gls{uav}s \cite{morge-rollet_drone_2022, li_dim_2025}, y constituye la base del detector implementado en esta investigación. Los parámetros que determinan el funcionamiento de este bloque funcional no se ajustan de forma automática, ya que esto supondría una fase de entrenamiento adicional que arrastraría todos los problemas comentados en secciones anteriores, en lugar de ser una herramienta de apoyo al modelo implementado. En este contexto, se trata de un algoritmo determinista, cuyas constantes de operación se derivan analíticamente del estudio realizado sobre el conjunto de datos seleccionado, de las estadísticas del ruido gaussiano presente en él y de la geometría espectral de las distintas transmisiones \gls{fhss} identificadas.

\subsection{Fundamento físico: entropía como medida de concentración espectral}
\label{subsec:fundamento_entropia}

Considérese la representación tiempo-frecuencia obtenida mediante la {\gls{stft}}. En cada ventana temporal $m$, se extrae la densidad espectral de potencia $\mathbf{P}_m = \{P_{k,m}\}_{k=0}^{N_w-1}$, compuesta por $N_w$ \textit{bins}. La fracción de potencia normalizada por \textit{bin} se define como:

\begin{equation}
    p_{k,m} = \frac{P_{k,m}}{\displaystyle\sum_{k=0}^{N_w-1} P_{k,m}}, \quad \forall k \in \{0, \ldots, N_w-1\}
    \label{eq:prob_espectral}
\end{equation}

Dado que la suma de las fracciones $p_{k,m}$ a lo largo de los $N_w$ \textit{bins} es exactamente 1, el vector resultante opera matemáticamente como una función de masa de probabilidad. Esto permite aplicar métricas de teoría de la información clásica directamente al dominio frecuencial de la señal \cite{morge-rollet_drone_2022}. En concreto, la entropía espectral cuantifica la uniformidad con la que se distribuye la energía en la banda analizada, y se calcula en bits por ventana temporal como:

\begin{equation}
    H(m) = -\sum_{k=0}^{N_w-1} p_{k,m} \cdot \log_2\!\left(p_{k,m}\right)
    \label{eq:entropia_shannon}
\end{equation}

Esta magnitud posee dos propiedades analíticas directamente explotables para la detección:

\begin{itemize}\itemsep 0em
    \item Máximo de entropía (ruido blanco): cuando la energía se distribuye de forma uniforme en todos los \textit{bins} (distribución espectral plana, característica del ruido blanco gaussiano), la distribución $\{p_{k,m}\}$ converge a la distribución uniforme y la entropía alcanza su cota superior teórica: $H_{\max} = \log_2(N_w)$ bits \cite{cover_elements_2006}. Para $N_w = 2048$ \textit{bins}, $H_{\max} = 11$ bits.

    \item Mínimo de entropía (transmisión activa): durante una transmisión \gls{fhss}, la energía se concentra en un conjunto reducido de \textit{bins} de la portadora activa, concentrando la distribución de potencia en unas frecuencias determinadas. En ese instante, $H(m) \ll H_{\max}$, generando un mínimo local en la curva $\{H(m)\}$ que señala inequívocamente la presencia de una ráfaga. Esta propiedad ha sido formalmente caracterizada como test de hipótesis para la detección de señales correladas en presencia de \gls{awgn} \cite{zhang_entropy-based_1989}.
\end{itemize}

La idea intuitiva es directa: el ruido térmico distribuye su potencia en todos los \textit{bins} simultáneamente (máxima incertidumbre), mientras que la transmisión de un dispositivo que use \gls{fhss} concentra la información en un conjunto reducido de frecuencias (mínima incertidumbre). El detector explota esta asimetría sin requerir conocimiento previo alguno sobre el protocolo, la frecuencia portadora ni la modulación empleada por el transmisor.

\subsubsection{Blanqueo espectral: robustez frente a la entropía estándar}
\label{subsubsec:whitening}

La estimación directa de la entropía sobre el periodograma empírico obtenido mediante la \gls{stft} adolece de una degradación sistemática originada por la falta de planicidad espectral del canal de comunicaciones. Fenómenos como el desvanecimiento selectivo en frecuencia impuesto por el medio de propagación, las tolerancias de los osciladores y la respuesta en frecuencia no ideal de las distintas cadenas de filtros del receptor \gls{sdr} introducen alteraciones energéticas localizadas en intervalos espectrales específicos \cite{tandra_snr_2008}. Estas variaciones determinísticas inducen una entropía base inferior a la cota teórica de máxima incertidumbre $H_{\max}$, incrementando la varianza temporal del estimador y comprometiendo la detectabilidad de las ráfagas transmitidas en regímenes de baja \gls{snr}.

La solución implementada es el blanqueo espectral: antes de calcular la entropía, cada \textit{bin} se normaliza por su mediana\footnote{En estadística, la mediana muestral se prefiere frente a la media aritmética por su elevado punto de ruptura (se define el punto de ruptura o de quiebre de un estimador como la proporción de observaciones incorrectas que puede manejar antes de proporcionar un resultado incorrecto) del 50\,\%. En nuestro escenario, las transmisiones de alta energía (las transmisiones \gls{fhss} o ráfagas WiFi) actúan como valores atípicos (conocidos en el ámbito del \textit{Machine Learning} como \textit{outliers}) temporales. Con el uso de la media, estas ráfagas elevarían artificialmente el nivel de referencia del \textit{bin}, sesgando la normalización y atenuando el contraste de las ráfagas más débiles. La mediana permanece inalterada frente a estos eventos, estimando de forma insesgada el verdadero suelo de ruido siempre que la ocupación del canal no supere la mitad de la ventana temporal.} de potencia a lo largo del eje temporal:

\begin{equation}
    \tilde{P}_{k,m} = \frac{P_{k,m}}{\tilde{P}_{k}^{\text{med}}}, \qquad \tilde{P}_{k}^{\text{med}} = \operatorname{mediana}_{m}\!\left(P_{k,m}\right)
    \label{eq:whitening}
\end{equation}

Esta operación normaliza el ruido blanco a una potencia espectral unitaria en todos los \textit{bins}, independientemente de la forma espectral del canal. En regímenes de \gls{snr} favorables o moderadamente adversos, la aparición de una ráfaga \gls{fhss} provoca que los \textit{bins} de la portadora activa superen de manera clara el nivel de referencia blanqueado, amplificando el contraste de la caída de entropía frente a la línea de base. Sin embargo, a medida que la señal se degrada hacia escenarios de \gls{snr} muy desfavorables (en torno a $-12$\,dB), la energía de la ráfaga queda progresivamente enmascarada por la varianza del ruido térmico, provocando que el decremento de entropía sea sutil o marginal, pero notable. El vector de entropía blanqueada, empleado en el resto del \textit{pipeline}, se calcula sustituyendo $P_{k,m}$ por $\tilde{P}_{k,m}$ en las Ecuaciones~\ref{eq:prob_espectral} y~\ref{eq:entropia_shannon}.

\subsection{Parámetros de la \gls{stft} y Resolución Tiempo-Frecuencia}
\label{subsec:parametros_stft}
La \gls{stft} se evalúa sobre la señal recibida en banda base $r[n] = r_I[n] + j\,r_Q[n]$. El parámetro de diseño fundamental es la longitud de la ventana de análisis $N_w$, que rige el compromiso entre la resolución temporal y frecuencial \cite{oppenheim_discrete_2014}:

\begin{equation}
    \Delta f = \frac{F_s}{N_w}, \qquad \Delta t \approx \frac{N_w}{2 F_s}
    \label{eq:resolucion_stft}
\end{equation}

Para el conjunto de datos \textit{NoisyUAV v2}, con datos muestreados a $F_s = 14$\,MHz, se propone una ventana de $N_w = 2048$ con un solapamiento del 50\,\% ($H = 1024$). Esta configuración proporciona una resolución frecuencial de $\Delta f \approx 6.84$\,kHz/\textit{bin} y temporal de $\Delta t \approx 0.073$\,ms/\textit{frame}. El solapamiento es un factor de diseño crítico para asegurar la detección en el peor escenario posible: la ráfaga \gls{fhss} más corta registrada en el conjunto de datos (0.44\,ms) abarca 6 \textit{frames} consecutivos. Esto garantiza que cualquier transmisión activa quede muestreada íntegramente, con independencia de su alineación respecto al inicio de la ventana de análisis. Simultáneamente, la resolución $\Delta f$ permite discriminar la energía de la portadora, cuyo ancho de banda en protocolos de radiocontrol abarca varios megahercios \cite{morge-rollet_drone_2022}.

\vspace{\baselineskip}

Para extraer automáticamente las ráfagas representadas en la Figura \ref{fig:ed_vs_entropy_cfar}, la señal procesada no se evalúa directamente, sino que atraviesa una cadena de preprocesado que se describe a continuación. En primer lugar, se aplica un suavizado sobre la curva para atenuar la varianza de alta frecuencia del ruido térmico. Posteriormente, el sistema estima el suelo de ruido local ($\widehat{H}_{90}$ en la leyenda) para establecer la línea base del canal. A partir de este nivel, se calcula un umbral adaptativo (\gls{cfar}) ($au_H$) que mantiene constante la probabilidad de falsa alarma ante variaciones del entorno. Finalmente, las detecciones que cruzan el umbral atraviesan un filtro de duración mínima diseñado para ignorar el ruido impulsivo; esto explica por qué ciertas disminuciones de entropía visibles en la gráfica no activan el detector. La formulación matemática de estas etapas y la definición paramétrica de los filtros se desarrollan de forma exhaustiva en la Sección~\ref{sec:preprocesamiento_analitico} del Capítulo~\ref{cap:arquitectura}.

Para ilustrar la superioridad del detector implementado en este trabajo, se establece una comparativa frente a un detector de energía clásico. Con el objetivo de garantizar la validez de la evaluación, este modelo de referencia opera bajo la misma arquitectura de decisión: evalúa la energía suavizada $P_s(m)$ y aplica parámetros idénticos para la estimación del nivel base de ruido, el umbral adaptativo \gls{cfar} y el filtro de duración mínima. A pesar de operar en igualdad de condiciones, el detector de energía (panel central) es incapaz de extraer las transmisiones de los drones debido a la presencia de interferencias y un nivel muy bajo de \gls{snr} ($-14~dB$). En este escenario, la potencia del dron se encuentra físicamente por debajo de la línea base del ruido del canal, produciendo que el detector clásico solo extraiga las emisiones de alta potencia (Bluetooth y Wi-Fi).

Por el contrario, el detector de entropía espectral propuesto (panel inferior) localiza con precisión las cuatro transmisiones del dron y la ráfaga Bluetooth. De forma simultánea, ignora la emisión Wi-Fi: al utilizar modulación OFDM, su ancho de banda genera una distribución espectral plana muy similar a la del ruido (la potencia se distribuye equitativamente en todo el ancho de banda analizado, no se concentra en un conjunto de frecuencias en específico), lo que impide una caída de entropía significativa.

\begin{figure}[!htbp]
    \centering
    \includegraphics[width=0.88\linewidth]{figuras/seccion3/ed_vs_entropy_cfar_edit.pdf}
    \caption{Comparativa entre el detector de energía y el detector de entropía espectral propuesto ante un entorno de baja relación señal a ruido y con coexistencia de interferencias de banda estrecha (Bluetooth) y banda ancha (Wi-Fi).}
    \label{fig:ed_vs_entropy_cfar}
\end{figure}

\section{Redes Neuronales Convolucionales de Valores Complejos (CV-CNN)}
\label{sec:cvcnn}

El procesamiento de la señal I/Q con redes neuronales convencionales plantea un problema estructural. La señal recibida en banda base es una secuencia compleja $r[n] = r_I[n] + j\,r_Q[n]$, donde la magnitud $|r[n]|$ y la fase $\angle r[n]$ son interdependientes, ya que los protocolos de modulación digital codifican la información mediante rotaciones y escalados en el plano complejo. Procesar las componentes $r_I$ y $r_Q$ como si fueran dos canales reales independientes mediante convoluciones estándar rompe esta interdependencia geométrica. Como consecuencia, la red pierde la capacidad estructural para detectar rotaciones de fase y se ve forzada a reaprender unas correlaciones que la modulación ya incluye por diseño.

Para resolver este problema de representación, se emplean \glspl{cv-cnn} \cite{bassey_survey_2021}, en las que tanto las activaciones como los pesos de los filtros operan en el cuerpo de los números complejos ($\mathbb{C}$). Dado un tensor de entrada $z = x_R + jx_I$ y un filtro $W = W_R + jW_I$, la convolución compleja se computa directamente aplicando la propiedad distributiva de la multiplicación:

\begin{equation}
    W \ast z = (W_R \ast x_R - W_I \ast x_I) + j(W_R \ast x_I + W_I \ast x_R)
    \label{eq:convolucion_compleja}
\end{equation}

Esta operación combina cuatro convoluciones reales (denotadas por el operador $\ast$) que generan un acoplamiento cruzado entre las componentes en fase y en cuadratura. Este entrelazamiento obliga al filtro $W$ a comportarse como un operador de rotación y escalado bidimensional \cite{hirose_complex-valued_2012}. Mediante esta formulación, la red preserva e impone desde su propia arquitectura las relaciones geométricas inherentes a las transmisiones de \gls{rf}.

La formulación compleja aporta una ventaja estructural en términos de eficiencia paramétrica. Una capa convolucional real que procesa $C_{\text{in}}$ mapas de características de entrada para producir $C_{\text{out}}$ mapas de salida mediante filtros de tamaño $K$ requiere $C_{\text{in}} \cdot C_{\text{out}} \cdot K$ parámetros. Para procesar las ramas $r_I$ y $r_Q$ por separado, una red real requiere duplicar ambas dimensiones, elevando el coste a $4 \cdot C_{\text{in}} \cdot C_{\text{out}} \cdot K$ parámetros. En contraste, la capa compleja equivalente está definida por pesos $W \in \mathbb{C}$ parametrizados por dos escalares reales $(W_R, W_I)$, por lo que opera con $2 \cdot C_{\text{in}} \cdot C_{\text{out}} \cdot K$ grados de libertad \cite{hirose_complex-valued_2012}. Al reducir el número de parámetros a la mitad frente a su equivalente real, se restringe estructuralmente la capacidad de la red para memorizar el ruido térmico. Este efecto regularizador intrínseco mejora la capacidad de generalización del modelo en escenarios de \gls{snr} negativa \cite{trabelsi_deep_2018, bassey_survey_2021}.

\subsection{Retropropagación mediante el Cálculo de Wirtinger}
El entrenamiento de la red requiere optimizar la función de pérdida $\mathcal{L}$ respecto a los pesos complejos $W$. Dado que $\mathcal{L}$ cuantifica un error, su salida es estrictamente real ($\mathcal{L}: \mathbb{C}^n \rightarrow \mathbb{R}$). Según las ecuaciones de Cauchy-Riemann, una función real de variable compleja solo es derivable (holomorfa) si es constante. En consecuencia, no es posible aplicar la derivada compleja estándar para calcular el gradiente \cite{bouboulis_extension_2011}.

El Cálculo de Wirtinger evita esta restricción matemática tratando al peso $W$ y a su complejo conjugado $W^*$ como variables independientes. El teorema fundamental de este formalismo establece que, para funciones de coste de valor real, la dirección matemática de máximo crecimiento en el plano complejo viene dada exclusivamente por la derivada respecto al conjugado $W^*$. Esta es la razón técnica por la que el operador gradiente se define siempre con el asterisco de conjugación \cite{bouboulis_extension_2011}:

\begin{equation}
    \frac{\partial \mathcal{L}}{\partial W^*} = \frac{1}{2}\left(\frac{\partial \mathcal{L}}{\partial W_R} + j\,\frac{\partial \mathcal{L}}{\partial W_I}\right)
    \label{eq:wirtinger}
\end{equation}

Esta derivada se calcula en la práctica aplicando la retropropagación estándar sobre las componentes reales de forma separada. Para una transformación interna $f$, la regla de la cadena de Wirtinger propaga el gradiente hacia atrás sumando las contribuciones de ambas variables independientes:

\begin{equation}
    \frac{\partial \mathcal{L}}{\partial z} = \frac{\partial \mathcal{L}}{\partial f} \frac{\partial f}{\partial z} + \frac{\partial \mathcal{L}}{\partial f^*} \frac{\partial f^*}{\partial z}
    \label{eq:regla_cadena_wirtinger}
\end{equation}

Finalmente, la actualización de los tensores se ejecuta mediante el Descenso de Gradiente Complejo (\gls{cgd}), restando al peso actual la dirección de la derivada de Wirtinger multiplicada por la tasa de aprendizaje $\eta$ \cite{bouboulis_extension_2011}:

\begin{equation}
    W_{t+1} = W_t - \eta\,\frac{\partial \mathcal{L}}{\partial W^*}
    \label{eq:complex_gradient_descent}
\end{equation}

Este esquema es formalmente equivalente al descenso de gradiente ordinario en $\mathbb{R}^2$, con la propiedad adicional de que opera directamente sobre la representación fasorial sin necesidad de convertir los pesos a su forma real en ningún punto del proceso \cite{bouboulis_extension_2011}. La convergencia del entrenamiento queda garantizada bajo las mismas condiciones que en redes reales.

\subsection{Normalización y Funciones de Activación}
El entrenamiento estable de redes profundas exige controlar la distribución de las activaciones internas entre capas. La Normalización Compleja por Lotes (\gls{cbn}) modela cada canal como una distribución bivariante con covarianza entre las partes real e imaginaria \cite{trabelsi_deep_2018}. Sin embargo, calcular matrices de covarianza $2 \times 2$ por cada lote y canal hace inviable el entrenamiento sobre las grabaciones del conjunto \textit{NoisyUAV v2}. El modelo propuesto emplea una \gls{cbn} diagonal que normaliza las varianzas de la parte real y la imaginaria de forma independiente asumiendo ortogonalidad local. Esta relajación reduce el coste de procesamiento en un orden de magnitud y mantiene la estabilidad numérica de los gradientes.

Respecto a las funciones de activación, el teorema de Liouville\footnote{En análisis complejo, este teorema demuestra que toda función entera (holomorfa en todo el plano $\mathbb{C}$) que esté acotada debe ser constante. En el ámbito del aprendizaje profundo, esto implica que las funciones de activación tradicionales, al ser acotadas y no lineales, jamás podrán ser holomorfas y carecerán de derivada compleja directa \cite{hirose_complex-valued_2012}.} establece que la única función holomorfa y acotada en todo el plano complejo es la función constante \cite{trabelsi_deep_2018}. En consecuencia, no existe ninguna función de activación compleja que sea simultáneamente holomorfa y no lineal. La arquitectura emplea la función \gls{crelu}, que aplica la rectificación lineal estándar de forma separada a la componente en fase y en cuadratura:

\begin{equation}
    \operatorname{CReLU}(z) = \max(0,\operatorname{Re}(z)) + j\max(0,\operatorname{Im}(z))
    \label{eq:crelu}
\end{equation}

Las alternativas habituales (\textit{modReLU} o \textit{zReLU}) aplican umbrales sobre la magnitud o la fase de la señal. En entornos de baja \gls{snr}, cuando la amplitud se aproxima a cero debido al ruido, el cálculo de la fase se vuelve inestable, lo que provoca cambios bruscos en las funciones de activación que dependen del ángulo de la fase, dificultando la convergencia de la red. La función \gls{crelu} evita este problema procesando directamente y por separado las componentes en fase y en cuadratura, lo que garantiza gradientes estables \cite{trabelsi_deep_2018, bassey_survey_2021}.

Al final de la etapa de extracción de características, antes de la capa de clasificación lineal, el tensor complejo se proyecta al dominio real extrayendo el módulo:

\begin{equation*}
    |z| = \sqrt{\operatorname{Re}(z)^2 + \operatorname{Im}(z)^2}
\end{equation*}

Esta proyección garantiza la invarianza del modelo frente a la rotación absoluta de la señal recibida. Durante el procesamiento intermedio, la red opera en el dominio complejo para preservar y extraer la fase relativa (la relación geométrica entre muestras adyacentes), magnitud que codifica las transiciones de la señal \gls{rf}. Sin embargo, en un enlace no colaborativo, la distancia física y la desincronización de los osciladores locales (\gls{cfo}) introducen una rotación global aleatoria en la constelación que la red es capaz de eliminar a partir del cálculo del módulo en la última etapa. Este mecanismo evita que el clasificador asocie orientaciones estocásticas de la constelación con un tipo de dron específico, previniendo el sobreajuste y forzando a la red a clasificar únicamente por los patrones de fase relativa extraídos previamente \cite{hirose_complex-valued_2012}.

\section{Síntesis del Marco Teórico y Decisiones de Diseño}
\label{sec:sintesis_cap3}

Este capítulo fundamenta teóricamente las tres decisiones de diseño que fundamentan el sistema de detección propuesto:

\begin{enumerate}
    \item \textbf{Problema de detección con múltiples interferencias y muy baja \gls{snr}:} El modelo del sistema (Sección~\ref{sec:modelo_sistema}) demuestra matemáticamente que la señal recibida está degradada por el canal, múltiples interferencias cocanal y ruido térmico. Esta superposición de variables invalida los detectores de energía clásicos, diseñados bajo la asunción de canal plano y ruido gaussiano blanco.

    \item \textbf{Ruido de etiquetado y segmentación temporal:} El análisis de ocupación de la banda sobre las transmisiones (Sección~\ref{sec:dataset_noisyuav}) revela que las señales de los drones representan, como máximo, un 11,73\,\% de la duración de las grabaciones que conforman el conjunto de datos seleccionado. Procesar una grabación completa de 75\,ms de duración obligaría a la red a optimizar sus pesos tomando como entrada un tensor mayoritariamente compuesto por muestras de \gls{awgn} y etiquetado como Dron, generando ruido de etiquetas dependiente de la instancia (Sección~\ref{sec:label_noise}). Para solventar este problema, el sistema integra un detector basado en la entropía espectral de Shannon (Sección~\ref{sec:detector_entropia}) que aísla las transmisiones \gls{fhss}, de forma que los tensores de entrada de la red pasan a contener únicamente transmisiones de dron o interferencias de dispositivos que usen ráfagas con características similares a las de los drones. Se ha demostrado que para una $\text{SNR} = -14\,dB$ y presencia de interferencias, esta técnica localiza de forma precisa transmisiones de dron.

    \item \textbf{Uso de \gls{cv-cnn}:} El sistema propuesto emplea una arquitectura \gls{cv-cnn} (Sección~\ref{sec:cvcnn}) que procesa la secuencia IQ original en banda base como principal fuente de obtención de características debido a que mantienen la información relativa a la fase de la señal y reducen la carga computacional disminuyendo el número de parámetros. En este capítulo, se han detallado conceptos clave como la actualización de los pesos y las funciones de activación usadas en este tipo de redes.
\end{enumerate}